\subsection{Depth-bounded equality}\label{sec:visibility}

The preceding results describe full equational consequence. We now resolve
that consequence by the depth of terms allowed in a derivation. This small
filtration supplies the common example and notation for the companion
paper~\cite{ShcherbakovII}; no general depth theory is needed here.

Constants have depth zero and
$\depth f(t_1,\ldots,t_k)=1+\max_i\depth t_i$.
Let $T^{(D)}$ be the set of closed well-sorted terms of depth at most $D$.
On this finite set, $\equiv_D$ is the least equivalence relation containing
the ground equations whose sides belong to $T^{(D)}$, and closed under
operation applications whenever both resulting terms are in $T^{(D)}$.
For the carrier names $C_A$, put
\[
R_A^{(D)}=\equiv_D\big|_{C_A}.
\]
This is a bound on term depth within a derivation, not on its number of steps.

\begin{lemma}[Ground locality]\label{lem:locality}
Let $E$ be a finite ground equational presentation and $h(E)$ the maximum
depth of an axiom side (zero if $E$ is empty). If $s,t$ have depth at most
$d$ and $E\vdash s=t$, then $s\equiv_H t$ for
$H=\max\{d,h(E)\}$.
\end{lemma}

\begin{proof}
Let $V$ consist of all subterms of axiom sides and of $s,t$. Close the given
equations under congruence only at operation nodes represented in $V$.
The quotient of $V$ is a partial algebra: if two represented operation
nodes have equal argument classes, their result classes agree by closure.
For each sort, take these classes as a carrier, adjoining a default element
if no class exists. On a tuple represented by a node of $V$, use its result
class; on any other tuple, choose an arbitrary default of the result sort.
Interpret constants not in $V$ arbitrarily.
Structural induction shows that every term of $V$ evaluates to its class,
so this total algebra is a model of $E$. If the classes of $s$ and $t$ were
different, it would separate them, contradicting $E\vdash s=t$.
Thus their classes merge. All nodes used by this closure have depth at most
$H$, and each closure step is an allowed bounded-congruence step.
This proves the assertion and the usual ground subterm property
\cite{NelsonOppen1980,DowneySethiTarjan1980}.
\end{proof}

For each $b$, let $P_b=\{(a,c):\alpha(b,a)=c\in\Gamma\}$.
Define $\theta_{\rm val}$ to be the least native carrier congruence closed
under the rule
\begin{equation}\label{eq:valuecoherence}
a\equiv_\theta a',\quad (a,c),(a',c')\in P_b
\quad\Longrightarrow\quad c\equiv_\theta c'.
\end{equation}
It exists, for example by intersecting all congruences satisfying the rule.

\begin{theorem}[Three structural horizons]\label{thm:horizons}
In the typed pure-row fragment,
\[
R_A^{(0)}=\Delta_A,\qquad
R_A^{(1)}=\theta_{\rm val},\qquad
R_A^{(2)}=\theta_E^A.
\]
Horizon one imposes value coherence of supplied cells. Horizon two also
imposes homomorphic coherence and kernel feedback.
\end{theorem}

\begin{proof}
There are no equations between depth-zero carrier terms in the presentation,
so horizon zero preserves all carrier names. At horizon one, native ground
tables and supplied cells are available, but the compatibility equations
have sides of depth two. Congruence under native operations imposes native
carrier congruence. Congruence under $\alpha(b,-)$ imposes exactly
\eqref{eq:valuecoherence} on prescribed outputs. No further carrier equality
is necessary: interpret depth-one row nodes as separate copies of the
quotient by $\theta_{\rm val}$, identifying a row point with a carrier point
only where a supplied cell requires it. This finite partial algebra realizes
all horizon-one nodes without further carrier identifications.
Finally, every axiom side in the pure-row fragment has depth at most two,
so Lemma~\ref{lem:locality} gives the horizon-two assertion.
\end{proof}

\subsection{Exact quotients of the named interface}

Let
\[
\mathcal I=C_A\cup\{\alpha(b,a):b\in B,\ a\in A\},\qquad
X_D=\mathcal I/(\equiv_D|_{\mathcal I})\quad(D\ge1).
\]
We treat $X_D$ as a quotient of a labelled set. The interface need not be
closed under arbitrary native operations. Set
\[
\theta_1=\theta_{\rm val},\quad Q_1=A/\theta_1,
\qquad \theta_2=\theta_E^A,\quad Q_2=A/\theta_2.
\]
Let $U_b\subseteq Q_1$ consist of the classes containing at least one
supplied input for row $b$. Rule~\eqref{eq:valuecoherence} makes the supplied
output a well-defined map $v_b:U_b\to Q_1$; $U_b$ need not be a subalgebra.
Over $Q_2$, use the stabilized data $S_b,h_b,K_b,\kappa_b$ already defined.

\begin{theorem}[Named-interface structure and counts]\label{thm:interface}
The following are exact descriptions as labelled quotient sets.
\begin{enumerate}[label=(\roman*)]
\item $X_1$ consists of the central $Q_1$ and one copy of $Q_1$ for each
row, with the point $x\in U_b$ identified with $v_b(x)$ in the centre.
There are no other identifications. Consequently
\begin{equation}\label{eq:N1}
N_1:=|X_1|=|Q_1|+\sum_{b\in B}(|Q_1|-|U_b|).
\end{equation}
\item $X_2$ consists of the central $Q_2$ and the row copies
$Q_2/\kappa_b$, each identified with the centre along
$S_b/K_b\cong h_b(S_b)$. Distinct external row parts remain disjoint.
Consequently
\begin{equation}\label{eq:N2}
N_2:=|X_2|=|Q_2|+
\sum_{b\in B}\bigl(|Q_2/\kappa_b|-|S_b/K_b|\bigr).
\end{equation}
For an empty row the corresponding contribution is $|Q_2|$.
\item The relation on $\mathcal I$ is final at horizon two. Its first
stabilization horizon $\delta_E^{\mathcal I}\ge1$ satisfies
\[
\delta_E^{\mathcal I}=
\begin{cases}
1,&N_1=N_2,\\
2,&N_1>N_2.
\end{cases}
\]
\end{enumerate}
\end{theorem}

\begin{proof}
The horizon-one construction in Theorem~\ref{thm:horizons} gives (i).
Although different inputs can share a supplied output, each supplied-input
class removes precisely one otherwise new row point from the count.
For (ii), apply Theorem~\ref{thm:globalrows}. Each row has
$|Q_2/\kappa_b|$ image points, of which exactly $|S_b/K_b|$ are central;
external points from different rows do not overlap. Ground locality makes
this full initial equality already visible at horizon two. Finally,
$\equiv_1|_{\mathcal I}\subseteq\equiv_2|_{\mathcal I}$ on the same finite
labelled set. These relations are equal if and only if their numbers of
classes agree, which proves (iii).
\end{proof}

For a cell, define its first carrier-contact horizon by
\[
\chi_E(b,a)=\min\{D\ge1:\exists c\in A,\ 
\alpha(b,a)\equiv_D c\},
\]
with value $\infty$ when the set is empty. The theorem gives the exact rule
\[
\chi_E(b,a)=
\begin{cases}
1,&[a]_{\theta_1}\in U_b,\\
2,&[a]_{\theta_1}\notin U_b\text{ and }[a]_{\theta_2}\in S_b^{\kappa_b},\\
\infty,&\text{otherwise}.
\end{cases}
\]
For an empty row, only the last case applies. Similarly, for observed terms
$s,t$, their equality-visibility depth is
$\lambda_E(s,t)=\min\{D:s\equiv_D t\}$, or $\infty$ if they are unequal
in the initial algebra.

In Example~\ref{ex:cepobstruction}, the exact interface partitions are
\begin{center}
\begin{tabular}{@{}clc@{}}
\toprule
Horizon & Classes on $\{0,1,2,3,r_0,r_1,r_2,r_3\}$ & Count\\
\midrule
$1$ & $\{0,r_0,r_3\}\mid\{1\}\mid\{2,r_2\}\mid\{3\}\mid\{r_1\}$ & $5$\\
$2$ & $\{0,2,r_0,r_1,r_2,r_3\}\mid\{1\}\mid\{3\}$ & $3$\\
\bottomrule
\end{tabular}
\end{center}
Thus $\lambda_E(0,2)=2$, $\chi_E(b,1)=2$, and
$\delta_E^{\mathcal I}=2$. The two short derivations in that example use
only terms of depth at most two; the horizon-one partition proves that this
depth is necessary.

The feedback iteration count is a different quantity. The family in
Proposition~\ref{prop:sharpfeedback} requires $|A|-1$ strict simultaneous
quotient updates, although all its carrier equalities are visible at depth
two. A shallow derivation can contain many inference steps, and a structural
fixed-point computation can revisit the same shallow constraints.
